\begin{thebibliography}{99}

\bibitem[Acemoglu(2025)]{acemoglu2025}
Acemoglu, D. (2025).
\newblock The simple macroeconomics of AI.
\newblock \emph{Economic Policy}, 40(121), 13--58.

\bibitem[Acemoglu and Restrepo(2018)]{acemoglu2018}
Acemoglu, D. and Restrepo, P. (2018).
\newblock The race between man and machine: Implications of technology for growth, factor shares, and employment.
\newblock \emph{American Economic Review}, 108(6), 1488--1542.

\bibitem[Acemoglu and Restrepo(2022)]{acemoglu2022tasks}
Acemoglu, D. and Restrepo, P. (2022).
\newblock Tasks, automation, and the rise in US wage inequality.
\newblock \emph{Econometrica}, 90(5), 1973--2016.

\bibitem[Acemoglu et~al.(2022)]{acemoglu2022vacancies}
Acemoglu, D., Autor, D., Hazell, J. and Restrepo, P. (2022).
\newblock Artificial intelligence and jobs: Evidence from online vacancies.
\newblock \emph{Journal of Labor Economics}, 40(S1), S293--S340.

\bibitem[Agrawal et~al.(2018)]{agrawal2018}
Agrawal, A., Gans, J. and Goldfarb, A. (2018).
\newblock \emph{Prediction Machines: The Simple Economics of Artificial Intelligence}.
\newblock Harvard Business Review Press, Boston, MA.

\bibitem[Autor et~al.(2003)]{autor2003}
Autor, D.~H., Levy, F. and Murnane, R.~J. (2003).
\newblock The skill content of recent technological change: An empirical exploration.
\newblock \emph{The Quarterly Journal of Economics}, 118(4), 1279--1333.

\bibitem[Autor(2024)]{autor2024}
Autor, D. (2024).
\newblock Applying AI to rebuild middle class jobs.
\newblock NBER Working Paper No.~32140, National Bureau of Economic Research, Cambridge, MA.

\bibitem[Bastani and Waldenstr\"{o}m(2024)]{bastani2024}
Bastani, S. and Waldenstr\"{o}m, D. (2024).
\newblock AI, automation and taxation.
\newblock CESifo Working Paper No.~11084 / IZA Policy Paper No.~212, Munich and Bonn.

\bibitem[Briggs and Kodnani(2023)]{briggs2023}
Briggs, J. and Kodnani, D. (2023).
\newblock The potentially large effects of artificial intelligence on economic growth.
\newblock Goldman Sachs Global Economics Analyst, Goldman Sachs.

\bibitem[Brynjolfsson et~al.(2025a)]{brynjolfsson2025}
Brynjolfsson, E., Chandar, B. and Chen, R. (2025a).
\newblock Canaries in the coal mine? Six facts about the recent employment effects of artificial intelligence.
\newblock Working paper, Stanford Digital Economy Lab.

\bibitem[Brynjolfsson et~al.(2025b)]{brynjolfsson2025qje}
Brynjolfsson, E., Li, D. and Raymond, L. (2025b).
\newblock Generative AI at work.
\newblock \emph{The Quarterly Journal of Economics}, 140(2), 889--942.

\bibitem[Cazzaniga et~al.(2024)]{cazzaniga2024}
Cazzaniga, M., Jaumotte, F., Li, L., Melina, G., Panton, A.~J., Pizzinelli, C., Rockall, E. and Tavares, M.~M. (2024).
\newblock Gen-AI: Artificial intelligence and the future of work.
\newblock IMF Staff Discussion Note SDN/2024/001, International Monetary Fund, Washington, DC.

\bibitem[Doorley et~al.(2023)]{doorley2023automation}
Doorley, K., Gromadzki, J., Lewandowski, P., Tuda, D. and Van~Kerm, P. (2023).
\newblock Automation and income inequality in Europe.
\newblock IZA Discussion Paper No.~16499, Institute of Labor Economics, Bonn.

\bibitem[Doorley et~al.(2026)]{doorley2026}
Doorley, K., O'Connor, S., O'Shea, R. and Tuda, D. (2026).
\newblock Artificial intelligence and income inequality in Ireland.
\newblock Jointly-published Report 16, Economic and Social Research Institute and Department of Finance, Dublin.

\bibitem[DSIT and DfE(2023)]{dsit2023}
Department for Science, Innovation and Technology and Department for Education (2023).
\newblock The impact of AI on UK jobs and training.
\newblock UK Government, London.

\bibitem[Eloundou et~al.(2024)]{eloundou2023}
Eloundou, T., Manning, S., Mishkin, P. and Rock, D. (2024).
\newblock GPTs are GPTs: Labor market impact potential of LLMs.
\newblock \emph{Science}, 384(6702), 1306--1308.

\bibitem[Felten et~al.(2021)]{felten2021}
Felten, E., Raj, M. and Seamans, R. (2021).
\newblock Occupational, industry, and geographic exposure to artificial intelligence: A novel dataset and its potential uses.
\newblock \emph{Strategic Management Journal}, 42(12), 2195--2217.

\bibitem[Frey and Osborne(2017)]{frey2017}
Frey, C.~B. and Osborne, M.~A. (2017).
\newblock The future of employment: How susceptible are jobs to computerisation?
\newblock \emph{Technological Forecasting and Social Change}, 114, 254--280.

\bibitem[Goos and Manning(2007)]{goosmanning2007}
Goos, M. and Manning, A. (2007).
\newblock Lousy and lovely jobs: The rising polarization of work in Britain.
\newblock \emph{The Review of Economics and Statistics}, 89(1), 118--133.

\bibitem[Goos et~al.(2014)]{goos2014}
Goos, M., Manning, A. and Salomons, A. (2014).
\newblock Explaining job polarization: Routine-biased technological change and offshoring.
\newblock \emph{American Economic Review}, 104(8), 2509--2526.

\bibitem[Henseke et~al.(2025)]{henseke2025}
Henseke, G. et al. (2025).
\newblock How exposed are UK jobs to generative AI? The GAISI index.
\newblock arXiv:2507.22748.

\bibitem[Hosseini and Lichtinger(2026)]{hosseini2026}
Hosseini, S. and Lichtinger, G. (2026).
\newblock Generative AI as seniority-biased technological change: Evidence from U.S. r\'{e}sum\'{e} and job posting data.
\newblock SSRN Working Paper No.~5425555.

\bibitem[Humlum and Vestergaard(2025)]{humlum2025}
Humlum, A. and Vestergaard, E. (2025).
\newblock Large language models, small labor market effects.
\newblock Working paper, University of Chicago Becker Friedman Institute.

\bibitem[ILO(2025)]{ilo2025}
International Labour Organization (2025).
\newblock Generative AI and jobs: A 2025 update.
\newblock ILO, Geneva.

\bibitem[Klein~Teeselink(2025)]{kleinteeselink2025}
Klein~Teeselink, B. (2025).
\newblock Generative AI and labor market outcomes: Evidence from the UK.
\newblock SSRN Working Paper No.~5516798.

\bibitem[Korinek and Stiglitz(2019)]{korinek2019}
Korinek, A. and Stiglitz, J.~E. (2019).
\newblock Artificial intelligence and its implications for income distribution and unemployment.
\newblock In Agrawal, A., Gans, J. and Goldfarb, A. (eds), \emph{The Economics of Artificial Intelligence: An Agenda}, University of Chicago Press, Chicago, 349--390.

\bibitem[Moll et~al.(2022)]{moll2022}
Moll, B., Rachel, L. and Restrepo, P. (2022).
\newblock Uneven growth: Automation's impact on income and wealth inequality.
\newblock \emph{Econometrica}, 90(6), 2645--2683.

\bibitem[Noy and Zhang(2023)]{noy2023}
Noy, S. and Zhang, W. (2023).
\newblock Experimental evidence on the productivity effects of generative artificial intelligence.
\newblock \emph{Science}, 381(6654), 187--192.

\bibitem[OBR(2025)]{obr2025}
Office for Budget Responsibility (2025).
\newblock Fiscal risks and sustainability report.
\newblock OBR, London.

\bibitem[Pizzinelli et~al.(2023)]{pizzinelli2023}
Pizzinelli, C., Panton, A.~J., Tavares, M.~M., Cazzaniga, M. and Li, L. (2023).
\newblock Labor market exposure to AI: Cross-country differences and distributional implications.
\newblock IMF Working Paper WP/23/216, International Monetary Fund, Washington, DC.

\bibitem[Rockall et~al.(2025)]{rockall2025}
Rockall, E., Pizzinelli, C. and Tavares, M.~M. (2025).
\newblock Artificial intelligence and wage inequality in the United Kingdom.
\newblock IMF Working Paper, International Monetary Fund, Washington, DC.

\bibitem[Tolan et~al.(2021)]{tolan2021}
Tolan, S., Pesole, A., Mart\'{i}nez-Plumed, F., Fern\'{a}ndez-Mac\'{i}as, E., Hern\'{a}ndez-Orallo, J. and G\'{o}mez, E. (2021).
\newblock Measuring the occupational impact of AI: Tasks, cognitive abilities and AI benchmarks.
\newblock \emph{Journal of Artificial Intelligence Research}, 71, 191--236.

\bibitem[Webb(2019)]{webb2019}
Webb, M. (2019).
\newblock The impact of artificial intelligence on the labor market.
\newblock Working paper, Stanford University.

\bibitem[Williamson et~al.(2025)]{williamson2024}
Williamson, S., et~al. (2025).
\newblock Applying the complementarity-adjusted AI occupational exposure index to the United Kingdom and Ireland.
\newblock \emph{Journal for Labour Market Research}, 59, 30.

\end{thebibliography}
